\section{Estimated controls and centring: derivations}\label{app:controls}
The adjusted outcomes within a deletion unit share the error of the estimated control coefficient.
This appendix derives its effect on the fixed-offset estimator, states the additional term for the joint
point estimate and gives the approximations in the feasible variance calculation.
Here $z$ denotes the controls rather than the characteristics of the variance model.
\subsection{Setup}\label{ctl:setup}

Physical observation $\tau=1,\dots,F_i$ of deletion unit $i=1,\dots,n$ belongs to worker $g(i)$ and firm $j(i)$:
$y_{i\tau}=\alpha_{g(i)}+\psi_{j(i)}+z_{i\tau}'\gamma+\varepsilon_{i\tau}$, with $z_{i\tau}\in\mathbb{R}^k$. Errors have
mean zero and are independent across units; within a unit they may be dependent. Under match deletion a unit is a
mover's match ($F_i\ge1$) or one observation of a stayer, the rows of Section~\ref{sec:setting}; under observation
deletion every observation is a unit. Unit $i$ is one row,
\[
 x_i=\sqrt{F_i}\,(e_{g(i)}+e_{j(i)}),\qquad y_i(\gamma_0)=F_i^{-1/2}\sum_\tau\bigl(y_{i\tau}-z_{i\tau}'\gamma_0\bigr),\qquad
 \varepsilon_i=F_i^{-1/2}\sum_\tau\varepsilon_{i\tau},\qquad \sigma_i^2=\V[\varepsilon_i],
\]
with $\mu_i=x_i'\beta$, and with $S_{xx}$, $P=XS_{xx}^{-1}X'$, $M=I-P$, $B$ and the leave-out weights $C$ of
\eqref{eq:ustat} formed as in Section~\ref{sec:estimator}. For a known coefficient vector $\gamma_0$,
$\tilde\theta(\gamma_0)=y(\gamma_0)'Cy(\gamma_0)$ is the KSS estimator, unbiased at $\gamma_0=\gamma$.

$\hat\gamma$ is the coefficient of the full least-squares fit on worker dummies, firm dummies and the controls.
With $\tilde z_{i\tau}$ the controls after removing the two-way effects over physical observations and
$\Lambda=(\sum_{i,\tau}\tilde z_{i\tau}\tilde z_{i\tau}')^{-1}$, $\delta=\hat\gamma-\gamma=\Lambda\sum_{i,\tau}\tilde
z_{i\tau}\varepsilon_{i\tau}$. The estimator studied here is the fixed-offset one,
$\hat\theta=\tilde\theta(\hat\gamma)$.

Unit aggregates: $Z_i=F_i^{-1/2}\sum_\tau z_{i\tau}$, $\tilde Z_i=F_i^{-1/2}\sum_\tau\tilde z_{i\tau}$,
$\bar z_i=F_i^{-1}\sum_\tau z_{i\tau}$, $\omega_i=\sum_\tau(z_{i\tau}-\bar z_i)\varepsilon_{i\tau}$,
$G_i=\sum_\tau(z_{i\tau}-\bar z_i)(z_{i\tau}-\bar z_i)'$; $Z$ and $\tilde Z$ are the $n\times k$ matrices with rows
$Z_i'$ and $\tilde Z_i'$. Since the two-way part of $z$ is constant within a unit,
$\delta=\Lambda(\tilde Z'\varepsilon+\omega)$ with $\omega=\sum_i\omega_i$. Write $\varsigma_i=\Cov[\varepsilon_i,\omega_i]$,
$\Omega_i=\V[\omega_i]$, $K$ for the $n\times k$ matrix with rows $\varsigma_i'$, $\Omega=\sum_i\Omega_i$,
$V=\diag(\sigma_i^2)$ and $\Pi=I-\tilde Z\Lambda Z'$. For $F_i=1$, $\omega_i=0$.

\subsection{The variance of the fixed-offset estimator}\label{ctl:var}

\paragraph{The noise of the adjusted outcome.} $y(\hat\gamma)=\mu+v$ with
\begin{equation}\label{ctl:eq:v}
 v=\varepsilon-Z\delta=\Pi'\varepsilon-Z\Lambda\omega,\qquad
 S=\V[v]=\Pi'V\Pi-\Pi'K\Lambda Z'-Z\Lambda K'\Pi+Z\Lambda\Omega\Lambda Z'.
\end{equation}
$S$ equals $V$ plus a symmetric matrix of rank at most $2k$. The estimator is the quadratic form
$\hat\theta=(\mu+v)'C(\mu+v)$, so with $m=C\mu$
\begin{equation}\label{ctl:eq:exact}
 \hat\theta-\theta=2m'v+v'Cv,\qquad \E\hat\theta-\theta=\tr(CS),\qquad
 \V[\hat\theta]=4m'Sm+2\tr(CSCS)\ \text{ for Gaussian errors.}
\end{equation}
The bias $\tr(CS)$ can be nonzero because $S$ is not diagonal.

\paragraph{The linear part.} $2m'v=\sum_i(2m_i^*\varepsilon_i+c_m'\omega_i)$ with $m^*=\Pi m$, $c_m=-2\Lambda Z'm$, and
\begin{equation}\label{ctl:eq:Vlin}
 4m'Sm=\sum_i\bigl(4m_i^{*2}\sigma_i^2+4m_i^*\varsigma_i'c_m+c_m'\Omega_ic_m\bigr).
\end{equation}
With known controls it is $4\sum_im_i^2\sigma_i^2$, the first term of \eqref{eq:V}. (For two
targets $s$ and $t$ it is $4m_s'Sm_t$.) The linear part is exactly uncorrelated with $\varepsilon'C\varepsilon$,
without Gaussianity: since $C_{ii}=0$, each term of $\varepsilon'C\varepsilon$ involves two different units, and
multiplying it by the linear contribution of any unit leaves an independent factor with mean zero. This does not
establish zero covariance with $v'Cv$ for non-Gaussian errors.

\paragraph{The quadratic part.} $v'Cv=\varepsilon'C\varepsilon+R$ with $R=-2\varepsilon'CZ\delta+\delta'Z'CZ\delta$,
a quadratic form of rank at most $2k$ in the physical errors. Low rank does not make it small: writing
$V_0=4m'Sm+2\tr(CVCV)$, what is needed is $\V[R]=o(V_0)$ and $\E R=o(\sqrt{V_0})$. \emph{These hold under
stated conditions, not in general.} In the Gaussian case it suffices that $k$ is fixed, the covariance of the
physical errors has bounded operator norm, $\|Z\Lambda^{1/2}\|_\op$ is bounded and $\|C\|_\op^2/V_0\to0$; then
$\V[R]=O(k\|C\|_\op^2)$ and $|\E R|=O(k\|C\|_\op)$. The conditions can fail when a dominant quadratic direction is
aligned with the controls. (For example, two workers, one firm, $z$ equal to the worker indicator
plus a within-worker contrast; $\max_i\tilde Z_i'\Lambda\tilde Z_i=1/n$, yet $\V[\hat\theta]/(2\tr(CVCV))\to4$
and the bias is of the order of the standard deviation.) Formula \eqref{ctl:eq:exact} with $S$ does not need this
paragraph or its conditions; this is why the estimator below is built on $S$ rather than on $V$ plus a
first-order term.

\subsection{The jointly estimated version}\label{ctl:joint}

For observation units the joint estimator is $y'C^{\rm f}y$ with the kernel of the full design $[X,\mathbf Z]$.
$B^{\rm f}=\Pi B\Pi'$ exactly, so the $B$-part of its linear influence equals that of $m^*$.  The
leave-out corrections differ: with $D=\diag(B_{ii}/M_{ii})$ from \eqref{eq:ustat},
$M^{\rm f}=M-\tilde Z\Lambda\tilde Z'$ and $D^{\rm f}=\diag(B^{\rm f}_{ii}/M^{\rm f}_{ii})$,
\begin{equation}\label{ctl:eq:joint}
 m_{\rm joint}-m^*=-\tfrac12M^{\rm f}\bigl[(D^{\rm f}-D)\mu+D^{\rm f}Z\gamma\bigr].
\end{equation}
The second term can be non-negligible. It depends on $M^{\rm f}D^{\rm f}Z\gamma$, so it vanishes when $D^{\rm f}$ is
a scalar matrix even if the control index $Z\gamma$ varies. It arises because the joint estimator uses the raw
outcome in its correction, so it is not invariant to adding $Za$ to the outcome, while the fixed-offset estimator
is. \textbf{The two estimators are therefore not first-order equivalent in general}, and the variance below
describes the fixed-offset estimator. (For example, workers with two or four observations, a
control alternating $\pm1$, unit error variance: $\V[\hat\theta_{\rm joint}]/\V[\hat\theta_{\rm fixed}]\to1+4\gamma^2/45$.)
In the check of Appendix~\ref{app:ctlnum} the joint variance is 0.995 to 1.035 times the fixed-offset one. A term
for \eqref{ctl:eq:joint} is not part of this analysis, nor of \fevc; Section~\ref{sec:impl-mc} reports what
the package simulations show of it.

\subsection{Estimating it within the method}\label{ctl:est}

The method supplies variance estimates $\st_i^2$ for the units (for the units of a repeat-measurement set $U$,
the sample variance $q_U$ of their adjusted outcomes, as $q_g$ in \eqref{eq:q}; a fitted value elsewhere),
$\tilde V=\diag(\st_i^2)$, the influence estimates $\hat m=Cy(\hat\gamma)$, and the pseudo-noise draws of
Section~\ref{sec:computation}. A set $U$ consists of the rows of a stayer or, under observation deletion, of a
worker--firm cell with two or more observations (Section~\ref{sec:setting}), and $T_U$ is its number of rows. Let
$\tilde e_i=\hat e_i/M_{ii}$, with $\hat e_i$ the residual of unit $i$ in the unit design, be its leave-out residual,
so that $\hat\sigma_i^2=y_i(\hat\gamma)\,\tilde e_i$ is its leave-out estimate; let
$\hat\omega_i=\sum_\tau(z_{i\tau}-\bar z_i)\bigl(y_{i\tau}-\bar y_i-(z_{i\tau}-\bar z_i)'\hat\gamma\bigr)$, with $\bar
y_i$ the mean of the unit's $y_{i\tau}$, be the realised within-unit product; and let $\hat c_m=-2\Lambda Z'\hat m$,
$\hat m^*=\Pi\hat m$ and $\hat\Omega=\sum_i\hat\omega_i\hat\omega_i'$.

\paragraph{Linear part.} The plug-in for $4m'Sm$ with realised products for $\varsigma_i$ and $\Omega_i$:
\begin{equation}\label{ctl:eq:Lhat}
 \hat L=\sum_i\bigl(4\hat m_i^{*2}\st_i^2+4\hat m_i^*\tilde e_i\,\hat\omega_i'\hat c_m+(\hat\omega_i'\hat c_m)^2\bigr).
\end{equation}
With $S$ known, $\E[4\hat m'S\hat m]=4m'Sm+4\tr(CSCS)$ exactly, since $\V[\hat m]=CSC$.

\paragraph{Quadratic part.} Draw $\zeta\sim N(0,\tilde V)$ and independent standard normal variables $\xi_i$,
and set
\begin{equation}\label{ctl:eq:draw}
 v^*=\zeta-Z\Lambda\bigl(\tilde Z'\zeta+\textstyle\sum_i\hat\omega_i\xi_i\bigr),\qquad
 \V^*[v^*]=S_d=\Pi'\tilde V\Pi+Z\Lambda\hat\Omega\Lambda Z' .
\end{equation}
With $u=Cv^*$, $\hat Q$ is the average over draws of $u'S_du=\sum_i\st_i^2(\Pi u)_i^2+
(\Lambda Z'u)'\hat\Omega(\Lambda Z'u)$, an estimate of $\tr(CS_dCS_d)$. The estimate is
\begin{equation}\label{ctl:eq:Vhat}
 \hat\V[\hat\theta]=\hat L-2\hat Q+\text{pair term},
\end{equation}
where the pair term is the last term of \eqref{eq:Vhat}, the adjustment for the pairs of rows of a set. Without
controls this is \eqref{eq:Vhat} with the double sum simulated as in \eqref{eq:sim}. With $S$ known and Gaussian
errors, $4\hat m'S\hat m-2\tr(CSCS)$ is exactly unbiased for \eqref{ctl:eq:exact}.

\paragraph{What is approximate in the feasible version.} The following discrepancies remain; their aggregate
effect has not been bounded. The combined procedure is checked in one Gaussian design
(Appendix~\ref{app:ctlnum}). Bounded whole-unit leverage is weaker than the vanishing maximum influence usually
needed to replace the moments by these estimates. Every feasible claim is for Gaussian errors: with non-Gaussian
physical errors the transformed quadratic form can have nonzero diagonal blocks within units, so third- and
fourth-cumulant contributions to its variance are not identically zero.
\begin{enumerate}\itemsep1pt
\item $S_d$ leaves out the $K$ terms of $S$; they are kept in $\hat L$ through the realised cross products, and the
draw is not changed. This does not compensate. Let $S_0=S+H$ be the covariance of the draw without $K$ when the
covariance inputs are known, $H=\Pi'K\Lambda Z'+Z\Lambda K'\Pi$, of rank at most $2k$. Even with the linear part
$4\hat m'S\hat m$, the trace based on $S_0$ leaves the discrepancy
$2\{\tr(CSCS)-\tr(CS_0CS_0)\}=-4\tr(CSCH)-2\tr(CHCH)$, whose sign is not controlled. Under bounded $\|S\|_\op$ and
$\|H\|_\op$ it is $O(k\|C\|_\op^2)$, negligible for fixed $k$ when $\|C\|_\op^2/\V[\hat\theta]\to0$. It can be first
order with a dominant direction, even when the control information is diffuse: with $S=I$, $S_0=I+w_1w_1'$ for a
diffuse unit vector $w_1$ and a kernel approaching $\lambda_1w_1w_1'$, the two traces approach $\lambda_1^2$ and
$4\lambda_1^2$.
The method meets that case through Section~\ref{ctl:q1}: when one direction carries more than ten per cent of the
estimated variance, the interval is the $q=1$ interval (Section~\ref{sec:impl}), whose leading-score variance
$\Sigma_{bb}=w_1'Sw_1$ is computed with the $K$ terms; the quadratic part enters only through the remainder
weights $C^{\rm r}$, to which the bound applies with $C^{\rm r}$ in place of $C$. The row $\varsigma_i'$ of $K$ is zero
when the errors of unit $i$ have one variance and are uncorrelated within the unit, because the within-unit control
deviations sum to zero; the discrepancy needs within-unit heteroskedasticity that moves with the controls, or
serial correlation. An algebraic repair is available and not implemented: with $\hat K$ the matrix with rows
$\tilde e_i\hat\omega_i'$ and $\hat H=\Pi'\hat K\Lambda Z'+Z\Lambda\hat K'\Pi$, use
$\hat Q_{\rm full}=\hat Q-2\tr(CS_dC\hat H)+\tr(C\hat HC\hat H)$, whose conditional mean is
$\tr\{C(S_d-\hat H)C(S_d-\hat H)\}$, while $\hat L=4\hat m'(S_d-\hat H)\hat m$ exactly. The middle term is the average
of $u'\hat Hu$ over the existing draws; the last needs $2k$ further applications of $C$, whereas the control term as implemented needs no
further solve. (Under the default \code{nuisance(joint)} \fevc{} adds an offset pass of two solves per
point-estimate probe, 400 by default, to form the fixed-offset units; \code{nuisance(fixedoffset)} avoids it.) $S_d-\hat H$ need not be
positive semidefinite, so it cannot be drawn from.
\item $\hat\omega_i=\omega_i-G_i\delta$, not $\omega_i$: the realised product is shrunk by the fit, as a residual is.
In the extreme case of one match carrying all information on $\gamma$, $\hat\omega_i=0$ while $\Omega_i>0$.
\item $\E[\tilde e_i\hat\omega_i]=\varsigma_i-\sigma^2G_i\Lambda\tilde Z_i/M_{ii}$ under homoskedastic independent errors
with variance $\sigma^2$: the cross product is biased by a control-leverage term.
\item The sample variance of a set is computed from adjusted outcomes: under homoskedasticity
$\E q_U=\sigma^2\big[1-\tr\big(\Lambda\sum_{i\in U}(Z_i-\bar Z_U)(Z_i-\bar Z_U)'\big)/(T_U-1)\big]$, with $\bar Z_U$ the
mean of the $Z_i$ over $U$, and fitting $\gamma$ makes different sets dependent. The pair term is kept unchanged;
its exact justification (constant kernel entries inside a set, independent sets) no longer applies.
\item The leave-out estimates $\hat\sigma_i^2$ used to fit the model for units outside sets have mean
$(SM)_{ii}/M_{ii}$ rather than $\sigma_i^2$.
\item Products of the estimated $\hat m^*,\hat c_m$ with the realised $\hat\omega_i,\tilde e_i$ contain fourth-order
interactions that $\hat Q$ does not remove.
\item One sample variance per set. For independent adjusted outcomes $Y_i$, $i\in U$, with equal means and the
true controls removed, $\E q_U=\bar\sigma_U^2=T_U^{-1}\sum_{i\in U}\sigma_i^2$. Pooling is then unbiased for the
linear term $\sum_{i\in U}a_i\sigma_i^2$, $a_i=4(m_i^*)^2$, exactly when
$\sum_{i\in U}(a_i-\bar a_U)(\sigma_i^2-\bar\sigma_U^2)=0$; with controls the weights $a_i$ vary inside a set. Zero
linear correlation of the variances with a control is not enough, because the weights are quadratic:
control-index values $(-1,0,1)$ and variances $(2,1,2)$ are linearly uncorrelated, yet with weights equal to the
squared index the weighted sum is $4$ and pooling gives $10/3$. Exact cancellation in each set is stronger than
needed: the aggregate discrepancy must be negligible relative to the variance. The condition covers the linear
part only; the pair term with $q_U^2$ also uses one variance per set. With known controls, an alternative is available: for sets of three or more
units, $\hat v_i^\times=(T_U(Y_i-\bar Y_U)^2-q_U)/(T_U-2)$ is unbiased for $\sigma_i^2$ under these conditions without
Gaussianity, but it can be negative, so it cannot feed the draw; the pair term would need a new calculation; and
with fitted controls it is not exactly unbiased.
\end{enumerate}

\paragraph{Weak identification.} Where the control coefficient is poorly identified, or the sensitivity of the
estimate to it is itself noisy in a small sample, the second-order noise dominates and the feasible variance
estimate \eqref{ctl:eq:Vhat} can be non-positive. \fevc{} then reports the component without a normal interval: the
interval is withheld, or the $q=1$ interval is used. This was observed on a test fixture with a control whose
within-unit variation was small (amplitude 0.13); a well-identified control on the same data gives agreement with
a dense oracle to 0.004 per cent at 32{,}768 draws on the exact route.

\subsection{The $q=1$ pieces}\label{ctl:q1}

This subsection uses the decomposition of Section~\ref{sec:q1}. For a component with leading direction $w_1$
(unit-space weights $w_{i1}$ of unit length; a function of the design) and eigenvalue $\lambda_1$,
$\hat\theta=\lambda_1\big(\hat b_1^{\,2}-\hat\V[\hat b_1]\big)+\hat\theta_{\rm r}$ with $\hat b_1=w_1'y(\hat\gamma)$,
$\hat\V[\hat b_1]=\sum_iw_{i1}^2\,\hat\sigma_i^2$ and $\hat\theta_{\rm r}$ the leave-out form with the remainder
weights $C^{\rm r}$; this is an identity. With estimated controls $\hat b_1-b_1=w_1'v$, where $b_1=w_1'\mu$, so
\[
 \Sigma_{bb}=w_1'Sw_1=\sum_i\bigl((\Pi w_1)_i^2\sigma_i^2+2(\Pi w_1)_i\varsigma_i'c_w+c_w'\Omega_ic_w\bigr),\qquad
 c_w=-\Lambda Z'w_1,
\]
exactly, and $\hat\theta_{\rm r}-\theta_{\rm r}=2(m^{\rm r})'v+v'C^{\rm r}v$ gives
$\Cov[\hat b_1,2(m^{\rm r})'v]=2w_1'Sm^{\rm r}$ and, for Gaussian errors,
$\V[\hat\theta_{\rm r}]=4(m^{\rm r})'Sm^{\rm r}+2\tr(C^{\rm r}SC^{\rm r}S)$, $\Cov[\hat b_1,v'C^{\rm r}v]=0$ and
$\E\hat\theta_{\rm r}-\theta_{\rm r}=\tr(C^{\rm r}S)$. Estimates: \eqref{ctl:eq:Lhat} and \eqref{ctl:eq:Vhat} with
$C^{\rm r}$ for $\Sigma_{\rm rr}$ (the same draws); for $\Sigma_{bb}$, as in \eqref{eq:q1est}, the larger of the fitted
version $\sum_i\big((\Pi w_1)_i^2\st_i^2+2(\Pi w_1)_i\tilde e_i\hat\omega_i'c_w+(\hat\omega_i'c_w)^2\big)$ and the version
with the leave-out estimates $\hat\sigma_i^2$ in place of $\st_i^2$; for $\Sigma_{b{\rm r}}$ the plug-in for
$2w_1'Sm^{\rm r}$. A positive definiteness check of the $2\times2$ matrix remains necessary: with realised products a
fitted $\Sigma_{bb}$ need not be positive.

\paragraph{The centre of the interval.} $\E\hat\theta-\theta=\lambda_1\{\Sigma_{bb}-\E\hat\V[\hat b_1]\}+\tr(C^{\rm r}S)$.
With controls the first term need not be small on the weak-direction scale:
$\E\hat\V[\hat b_1]=\sum_iw_{i1}^2(SM)_{ii}/M_{ii}$, which in general differs from $w_1'Sw_1$. This is a bias of the
\emph{point estimate}. The interval of KSS (their Section~7.1; Section~\ref{sec:q1int}) is built from
$(\hat b_1,\hat\theta_{\rm r})$ and the map $\lambda_1b^2+\theta_{\rm r}$; $\hat b_1$ is unbiased for $b_1$ and
$\hat\theta_{\rm r}$ for $\theta_{\rm r}$ up to $\tr(C^{\rm r}S)$, so the interval need not inherit the first term, and
the decomposition is left unchanged. What it needs, and what is \emph{not proved} here, is joint normality of
$(\hat b_1,\hat\theta_{\rm r})$ with the control parts included (the Lindeberg condition is now on
$(\Pi w_1)_i\varepsilon_i+c_w'\omega_i$, not on $w_{i1}$ alone), consistency of the three estimated entries, and a
centring condition for the remainder: with $\rho$ the correlation of $\hat b_1$ and $\hat\theta_{\rm r}$,
$|\tr(C^{\rm r}S)|/\sqrt{\Sigma_{\rm rr}(1-\rho^2)}\to0$ (if $|\rho|$ stays away from one,
$\tr(C^{\rm r}S)=o(\Sigma_{\rm rr}^{1/2})$ suffices), or that bias estimated and subtracted. Joint normality and a
consistent covariance do not alone give coverage when the expectation is displaced from the target. The interval
is built on the two-dimensional statistic $(\hat b_1,\hat\theta_{\rm r})$ and need not be centred on the reported
scalar estimate, which may keep a bias $\lambda_1\{\Sigma_{bb}-\E\hat\V[\hat b_1]\}$; there is no inconsistency in that.

\subsection{Centring}\label{sec:centring}

\fevc{} subtracts the mean wage before forming the bias correction (under \code{centering(mean)}, its default).
Subtracting a constant $c$ from every wage leaves $\hat\theta$ unbiased: with $h$ the vector with entries
$\sqrt{F_i}$, the exact weights of the four components satisfy $h'C\mu=0$ and $h'Ch=0$, so
$\E[(y-ch)'C(y-ch)]=\mu'C\mu=\theta$. All formulas of this note then hold with $y-ch$ in place of $y$, and with $m$
computed from the centred means. \fevc{} uses the sample mean $\bar y$ for $c$, and the baseline propositions treat the centring constant as fixed. Sampling the outcome also changes its mean. Let $c_0=h'\mu/n_{\rm py}$
be the mean that $\bar y$ estimates. Centring at $\bar y$ instead of $c_0$ changes the estimate by
$-2\,(\bar y-c_0)\,h'C\varepsilon$, whose mean is $-(2/n_{\rm py})\,h'VCh$ with $V=\diag(\sigma_i^2)$. When
the error variances are bounded and bounded away from zero this change is of the order of $\lambda_{\max}$, and
therefore negligible relative to the standard error under condition (b) of Proposition~\ref{prop:conc}.

The standard error of \fevc{} does not hold the mean fixed. Under \code{centering(mean)} it centres the
pseudo-noise of Section~\ref{sec:computation} in the same way as the outcome before it enters the kernel, so the
kernel is that of the centred estimate and no approximation that holds the observed mean fixed is made
(\code{e(inference\_mean\_omitted)} is zero). Under \code{centering(corrected)} the standard error is that of the
mean-centred estimate and treats the corrected-centring increment as fixed; \code{centering(none)} uses the
uncentred kernel. In the comparison of Section~\ref{sec:kss}, on two-period samples under observation deletion,
the centring leaves the estimate unchanged but not its standard error (Section~\ref{kss:sec:emp-se}).
